\chapter{Register Model}
\label{ch:registers}

\section{Overview}

The SIRC-1 CPU provides sixteen 16-bit registers, classified into three categories:

\begin{itemize}
	\item General-purpose registers (\reg{r1}--\reg{r7})
	\item Address register pairs (\reg{l}, \reg{a}, \reg{s}, \reg{p})
	\item Status register (\reg{sr})
\end{itemize}

All programmer-visible registers are 16 bits wide. Address registers can be used in pairs to form 24-bit addresses for
memory operations.

Only \reg{sr} (cleared to zero) and the \reg{p} pair (loaded from the reset vector) have defined values after reset;
all other registers are architecturally undefined. See Chapter~\ref{ch:exceptions}.

\section{Programming Model}

Figure~\ref{fig:user-programming-model} shows every register a protected-mode program can access directly: the seven
general-purpose registers, the low (unprivileged) word of each of the four address register pairs, and the
unprivileged low byte of \reg{sr}. \reg{sr}'s high byte is shown shaded because it reads as zero and cannot be
written directly in protected mode; see Section~\ref{sec:privilege-restrictions} below.

\begin{figure}[H]
	\centering
	\begin{bytefield}[bitwidth=1.1em]{16}
		\bitheader{0,15} \\
		\bitbox{16}{\reg{r1}} \\
		\bitbox{16}{\reg{r2}} \\
		\bitbox{16}{\reg{r3}} \\
		\bitbox{16}{\reg{r4}} \\
		\bitbox{16}{\reg{r5}} \\
		\bitbox{16}{\reg{r6}} \\
		\bitbox{16}{\reg{r7}} \\
		\bitbox{16}{\reg{ll} (link, low word)} \\
		\bitbox{16}{\reg{al} (address, low word)} \\
		\bitbox{16}{\reg{sl} (stack, low word)} \\
		\bitbox{16}{\reg{pl} (program counter, low word)} \\
		\bitbox{8}{\color{sircgray}\reg{sr} high byte (hidden)} & \bitbox{8}{\reg{sr} low byte}
	\end{bytefield}
	\caption{User Programming Model}
	\label{fig:user-programming-model}
\end{figure}

Figure~\ref{fig:supervisor-programming-model} shows what supervisor mode adds on top of the user programming model:
the high (privileged) word of each address register pair, and \reg{sr}'s high byte, both hidden from a protected-mode
program.

\begin{figure}[H]
	\centering
	\begin{bytefield}[bitwidth=1.1em]{16}
		\bitheader{0,15} \\
		\bitbox{16}{\reg{lh} (link, high word)} \\
		\bitbox{16}{\reg{ah} (address, high word)} \\
		\bitbox{16}{\reg{sh} (stack, high word)} \\
		\bitbox{16}{\reg{ph} (program counter, high word)} \\
		\bitbox{8}{\reg{sr} high byte} & \bitbox{8}{\color{sircgray}\reg{sr} low byte (same as user model)}
	\end{bytefield}
	\caption{Supervisor Programming Model Supplement}
	\label{fig:supervisor-programming-model}
\end{figure}

Table~\ref{tab:register-encoding} below lists every register's encoding and states exactly which ones are
privileged; see Section~\ref{sec:privilege-restrictions} for the full read/write privilege rules.

\section{General-Purpose Registers}

\subsection{Register Set}

The SIRC-1 provides seven general-purpose registers:

\begin{center}
	\begin{tabular}{cl}
		\toprule
		\textbf{Register} & \textbf{Description}       \\
		\midrule
		\reg{r1}          & General-purpose register 1 \\
		\reg{r2}          & General-purpose register 2 \\
		\reg{r3}          & General-purpose register 3 \\
		\reg{r4}          & General-purpose register 4 \\
		\reg{r5}          & General-purpose register 5 \\
		\reg{r6}          & General-purpose register 6 \\
		\reg{r7}          & General-purpose register 7 \\
		\bottomrule
	\end{tabular}
\end{center}

The programmer may use these registers for any purpose. They hold 16-bit values and may be used as source or
destination operands for all ALU instructions.

\subsection{Usage Conventions}

The CPU does not enforce any register-usage convention among \reg{r1}--\reg{r7}; all seven registers are equally
general purpose. Calling conventions and register-preservation rules are outside the scope of this document and are
addressed in the SIRC-1 ABI and Toolchain documentation.

\section{Address Register Pairs}

\subsection{Overview}

The SIRC-1 provides four address register pairs, each consisting of a high and low 16-bit register. When used together,
these pairs form 24-bit addresses by combining the low 8 bits of the high register with the full 16 bits of the low
register.

\begin{figure}[H]
	\centering
	\begin{bytefield}[bitwidth=0.8em]{32}
		\bitheader{0-31} \\
		\bitbox{8}{\color{sircgray}Unused} &
		\bitbox{8}{High Register} &
		\bitbox{16}{Low Register} \\
		\bitbox{8}{\color{sircgray}\tiny bits 31--24} &
		\bitbox{8}{\tiny bits 23--16} &
		\bitbox{16}{\tiny bits 15--0}
	\end{bytefield}
	\caption{Address Register Pair Formation}
	\label{fig:addr-pair}
\end{figure}

The high 8 bits (bits 31--24) of the high register are not used for addressing but remain available for storage,
enabling techniques such as tagged pointers.

The exact address composition and stored-address layout are specified in Chapter~\ref{ch:data-representation}.

Although each address register pair has a specific name and a suggested usage by convention, \mnemonic{STOR} and
\mnemonic{LOAD} instructions treat all four pairs identically, and any pair may be used for indirect addressing.

Loading a value directly into the \reg{p} pair performs a jump.

The link register pair updates automatically when \mnemonic{LDEL}, \mnemonic{LJSR}, or \mnemonic{BRSR} executes. The
program counter pair increments after every instruction.

\subsection{Link Register (l)}

\begin{center}
	\begin{tabular}{ll}
		\toprule
		\textbf{Component} & \textbf{Register}     \\
		\midrule
		High Register      & \reg{lh} (privileged) \\
		Low Register       & \reg{ll}              \\
		\bottomrule
	\end{tabular}
\end{center}

The link register pair stores return addresses for subroutine calls. When a subroutine is called using \mnemonic{LDEL},
\mnemonic{LJSR}, or \mnemonic{BRSR}, the return address is automatically saved to this register pair.

\textbf{Typical Use:}
\begin{itemize}
	\item Storing return addresses from subroutine calls
	\item Alternate storage for general memory addressing
\end{itemize}

\subsection{Address Register (a)}

\begin{center}
	\begin{tabular}{ll}
		\toprule
		\textbf{Component} & \textbf{Register}     \\
		\midrule
		High Register      & \reg{ah} (privileged) \\
		Low Register       & \reg{al}              \\
		\bottomrule
	\end{tabular}
\end{center}

This general-purpose address register pair is used for pointer operations and memory access.

\textbf{Typical Use:}
\begin{itemize}
	\item Base pointer for data structures
	\item Array indexing
	\item General memory addressing
\end{itemize}

\subsection{Stack Pointer (s)}

\begin{center}
	\begin{tabular}{ll}
		\toprule
		\textbf{Component} & \textbf{Register}     \\
		\midrule
		High Register      & \reg{sh} (privileged) \\
		Low Register       & \reg{sl}              \\
		\bottomrule
	\end{tabular}
\end{center}

The stack pointer register pair points to the current top of the stack. Pre-decrement and post-increment addressing
modes are particularly useful with this register.

\textbf{Typical Use:}
\begin{itemize}
	\item Stack operations (push/pop)
	\item Local variable allocation
	\item Function parameter passing
\end{itemize}

\subsection{Program Counter (p)}

\begin{center}
	\begin{tabular}{ll}
		\toprule
		\textbf{Component} & \textbf{Register}     \\
		\midrule
		High Register      & \reg{ph} (privileged) \\
		Low Register       & \reg{pl}              \\
		\bottomrule
	\end{tabular}
\end{center}

The program counter register pair points to the instruction currently being fetched. It is not modified during either
fetch phase; the low word \reg{pl} advances by 2 during the Decode and Register Fetch phase, and \reg{ph} is never
incremented, so a wrap stays inside the segment.

\textbf{Typical Use:}
\begin{itemize}
	\item Automatically updated during normal execution
	\item Modified by branch and jump instructions
	\item Used for position-independent code
\end{itemize}

\section{Status Register}

The status register (\reg{sr}) is a special 16-bit register that contains CPU status flags and control bits. It is
divided into two bytes:

\begin{itemize}
	\item \textbf{Low Byte (bits 7--0):} Condition flags, readable in both modes and updated as a side effect of ALU and
	      shift instructions in both modes; direct writes to \reg{sr} are privileged
	\item \textbf{High Byte (bits 15--8):} Control and privileged state
\end{itemize}

In protected mode, reads of \reg{sr} are redacted to the low byte (high byte reads as zero), and direct writes to
\reg{sr} raise a privilege violation fault. In supervisor mode, a direct write to \reg{sr} is permitted, except that it
cannot change the Exception Active (EA) bit.

The status register is discussed in detail in Chapter~\ref{ch:status-register}.

\section{Privilege Restrictions}
\label{sec:privilege-restrictions}

\subsection{High Address Register Restrictions}

The high component of each address register pair (\reg{lh}, \reg{ah}, \reg{sh}, \reg{ph}) is privileged:

\begin{itemize}
	\item In \textbf{supervisor mode}: Reading and writing are both allowed.
	\item In \textbf{protected mode}:
	      \begin{itemize}
		      \item Reading is allowed.
		      \item Direct writes trigger a privilege exception.
		      \item Address-register-pair write-back is allowed only when the high word would remain unchanged; otherwise it triggers a
		            privilege exception.
	      \end{itemize}
\end{itemize}

This restriction provides a form of memory segmentation and protection. User programs running in protected mode cannot
modify the segment portion of addresses, preventing them from accessing memory outside their designated segment.
Instructions such as \mnemonic{LDEA}, \mnemonic{LDEL}, \mnemonic{BRAN}, \mnemonic{BRSR}, \mnemonic{LJMP},
\mnemonic{LJSR}, and \mnemonic{RETS} remain usable in protected mode when they stay within the current segment. If
their address-register-pair write-back would alter a high word, the CPU raises a privilege violation fault.

The status register differs slightly from the high address registers:
\begin{itemize}
	\item High address registers (\reg{lh}, \reg{ah}, \reg{sh}, \reg{ph}) remain readable in protected mode, but direct writes
	      trigger a privilege violation fault.
	\item The status register's privileged byte is hidden on reads in protected mode, and direct writes in protected mode cannot
	      modify it.
\end{itemize}

\section{Register Addressing}

Registers are addressed using 4-bit register identifiers in instruction encoding:

\begin{table}[H]
	\centering
	\begin{tabular}{clc}
		\toprule
		\textbf{ID}  & \textbf{Register} & \textbf{Privileged}                                 \\
		\midrule
		\texttt{0x0} & \reg{r1}          & No                                                  \\
		\texttt{0x1} & \reg{r2}          & No                                                  \\
		\texttt{0x2} & \reg{r3}          & No                                                  \\
		\texttt{0x3} & \reg{r4}          & No                                                  \\
		\texttt{0x4} & \reg{r5}          & No                                                  \\
		\texttt{0x5} & \reg{r6}          & No                                                  \\
		\texttt{0x6} & \reg{r7}          & No                                                  \\
		\texttt{0x7} & \reg{sr}          & Reads: high byte redacted. Writes: fully privileged \\
		\texttt{0x8} & \reg{lh}          & Yes                                                 \\
		\texttt{0x9} & \reg{ll}          & No                                                  \\
		\texttt{0xA} & \reg{ah}          & Yes                                                 \\
		\texttt{0xB} & \reg{al}          & No                                                  \\
		\texttt{0xC} & \reg{sh}          & Yes                                                 \\
		\texttt{0xD} & \reg{sl}          & No                                                  \\
		\texttt{0xE} & \reg{ph}          & Yes                                                 \\
		\texttt{0xF} & \reg{pl}          & No                                                  \\
		\bottomrule
	\end{tabular}
	\caption{Register Encoding}
	\label{tab:register-encoding}
\end{table}

Because writes to \reg{sr} are fully privileged, any instruction whose destination field encodes \texttt{0x7} is
privileged in protected mode, even when the operation would only alter the unprivileged low byte of \reg{sr}. This
register-index assignment is deliberate: \reg{sr} does not occupy field \texttt{0x0} (unlike an earlier revision of
this encoding), so an all-zero instruction word decodes as \mnemonic{ADDI[N]} \reg{r1}, \texttt{\#0} -- a genuine,
always-unprivileged no-op. This is what lets erased or uninitialized memory (which reads back as all-zero words) be
executed safely instead of faulting; see the \mnemonic{NOOP} entry in Chapter~\ref{ch:meta-instructions}.

\section{Usage Examples}

\begin{example}
	\begin{lstlisting}
; Load immediate value into r1
LOAD r1, #100

; Add 2 to the value stored in r1 and store the result into r2
LOAD r2, r1
ADDI r2, #2

; Add r1 and r2, store in r3
ADDR r3, r1, r2
	\end{lstlisting}
	\caption{Basic Register Operations}
\end{example}

\begin{example}
	\begin{lstlisting}
; Load value from memory at address in 'a' register pair
LOAD r1, (a)

; Store r1 to memory at address (a + 8)
STOR (#8, a), r1

; Push r1 onto stack (pre-decrement)
STOR -(s), r1

; Pop stack into r2 (post-increment)
LOAD r2, (s)+

	\end{lstlisting}
	\caption{Address Register Usage}
\end{example}

\begin{example}
	\begin{lstlisting}
; Call subroutine at address in 'a' register
; Return address saved to 'l' register automatically
LJSR a

; ...subroutine code...

; Return from subroutine
; Copies 'l' register to 'p' register
RETS
	\end{lstlisting}
	\caption{Subroutine Call}
\end{example}
